\documentclass[a4paper]{article}

\usepackage{graphicx}
\usepackage{amsmath,amssymb}
\usepackage{minted}
\usepackage{multirow}
\usepackage{float}
\usepackage{algorithm}
\usepackage{algpseudocode}
\usepackage{todonotes}
\usepackage{forest}
\usepackage{xstring}
\usepackage{xcolor}
\usepackage[most]{tcolorbox}
\usepackage{xparse}
\usepackage{natbib}
\usepackage{adjustbox}
\usepackage[active,tightpage]{preview}

\PreviewBorder=4pt

\title{Pencil Contact Overlay Position in the iPad Image Editor}
\date{2026-10-13}

\begin{document}
   \maketitle

   \begin{preview}
      \section*{Fix}
      The raster Image editor previously passed Pencil touch coordinates
      directly from a centered or panned \texttt{UIScrollView} to a SwiftUI
      overlay covering the visible viewport. The scroll view has a
      nonzero \texttt{bounds.origin} when a small image, such as
      \(760 \times 500\) pixels, is centered on the iPad. The preview
      tool icon could consequently appear far above the actual Pencil tip.

      The owning Pencil quick-tool gesture controller now translates all
      Pencil touch and tap anchors from the scroll view's bounds coordinate
      system to visible viewport coordinates by subtracting
      \texttt{bounds.origin}. This normalization applies to first
      contact, double-tap palette placement, and dismissal taps, regardless
      of canvas dimensions, centering insets, zoom, or panning.

      Added a regression test for negative and positive scroll offsets
      and rebuilt the transfer archive as version 0.31.61 (build 88).
   \end{preview}
\end{document}
